\section*{Contoh}
\addcontentsline{toc}{section}{Contoh}

\begin{contoh}[Menokenisasi secara manual]
Diberikan program
\begin{lstlisting}[language=C]
int x = 10;
x = x + 1;
\end{lstlisting}
Token yang dihasilkan:
\begin{center}
\small
\begin{tabular}{@{}lll@{}}
\toprule
\textbf{Baris:Kolom} & \textbf{Token} & \textbf{Lexeme} \\
\midrule
1:1 & \code{TOK\_INT} & \code{int} \\
1:5 & \code{TOK\_IDENT} & \code{x} \\
1:7 & \code{TOK\_ASSIGN} & \code{=} \\
1:9 & \code{TOK\_INT\_LIT} & \code{10} \\
1:11 & \code{TOK\_SEMI} & \code{;} \\
2:1 & \code{TOK\_IDENT} & \code{x} \\
2:3 & \code{TOK\_ASSIGN} & \code{=} \\
2:5 & \code{TOK\_IDENT} & \code{x} \\
2:7 & \code{TOK\_PLUS} & \code{+} \\
2:9 & \code{TOK\_INT\_LIT} & \code{1} \\
2:10 & \code{TOK\_SEMI} & \code{;} \\
3:1 & \code{TOK\_EOF} & --- \\
\bottomrule
\end{tabular}
\end{center}
Perhatikan bahwa spasi tidak menghasilkan token, dan setiap token
memiliki lokasi awal.
\end{contoh}

\begin{contoh}[Token, lexeme, dan pattern]
Untuk token \code{FLOAT\_LIT}, lexeme \code{3.14} cocok dengan
pattern \verb![0-9]+\.[0-9]+!. Lexeme \code{42} bukan anggota token
ini karena tidak memiliki titik. Sebaliknya, \code{42} adalah lexeme
untuk token \code{INT\_LIT}. Pattern mendefinisikan himpunan lexeme;
token adalah nama himpunan tersebut.
\end{contoh}

\begin{contoh}[Kata kunci versus identifier]
Misalkan tabel kata kunci berisi \code{int}, \code{if}, dan
\code{while}. Masukan \code{int x;} menghasilkan token
\code{INT}, \code{IDENT(x)}, \code{SEMI}. Masukan \code{integer x;}
menghasilkan \code{IDENT(integer)}, \code{IDENT(x)}, \code{SEMI},
karena \code{integer} tidak ada dalam tabel kata kunci. Lexer tidak
menebak berdasarkan awalan; ia mencocokkan seluruh lexeme.
\end{contoh}

\begin{contoh}[Maximal munch]
Pada masukan \code{a<=b}, lexer membaca \code{<} lalu melihat
karakter berikutnya \code{=}. Karena \code{<=} cocok dan lebih
panjang daripada \code{<}, lexer memilih \code{LE}. Pada masukan
\code{a<<b}, tidak ada operator \code{<<} dalam Subset-C, sehingga
lexer memilih \code{LT} dua kali. Aturan maximal munch bekerja pada
pattern yang tersedia, bukan pada operator yang mungkin ada dalam
bahasa lain.
\end{contoh}

\begin{contoh}[Kesalahan leksikal]
Masukan \code{int a = 3.;} mengandung float tanpa digit setelah
titik. Lexer menghasilkan \code{TOK\_ERROR} pada posisi titik,
mencatat pesan \code{baris 1, kolom 10: digit diharapkan setelah
titik}, lalu melanjutkan. Parser menerima token kesalahan dan dapat
melaporkan masalah tanpa menghentikan seluruh kompilasi.
\end{contoh}
